\section{Research Gap Informed by the GLW Dataset and Literature}

Modern chronic myeloid leukemia (CML) guidelines and recent reviews emphasize response milestones, BCR--ABL1 molecular monitoring, deep molecular response (DMR), MR4.5, and treatment-free remission (TFR) as central treatment goals \cite{Hochhaus2020ELN,Jabbour2025JAMAReview,Haddad2026Algorithms,Senapati2023CMLManagement}. Early molecular response is also a clinically meaningful predictor of later outcomes and TFR eligibility \cite{Hughes2014EarlyMolecularResponse,Rasekh2026TFRPredictor,Ono2026MR45}. However, the GLW dataset shows that routine monitoring data have a more complex structure than the idealized response timelines often used in clinical summaries.

In the cleaned GLW cohort, 87 patients contributed 495 longitudinal molecular observations and 275 at-risk survival intervals. DMR was observed in 68 patients (78.2\%), while 19 patients were censored. Median follow-up was 39.0 months, and patients had a median of 5 visits. Monitoring was irregular: the median gap between visits was 6.0 months, 50.3\% of visit gaps exceeded 6 months, and 20.0\% exceeded 12 months. Therefore, first DMR is not observed as an exact event time for many patients; it is known only to occur within a monitoring interval.

This creates a data-driven research gap. Existing literature defines the clinical importance of DMR, MR4.5, CMR, and TFR, and recent work discusses assay performance and real-world monitoring \cite{Kim2026BCRABLqPCR,Verweij2026RemoteMonitoring,Parihar2026TFR,Atallah2021TFR}. Yet fewer applied CML studies directly model repeated log-MRD trajectories together with interval-observed DMR onset under irregular follow-up, assay-floor behavior, and incomplete baseline covariates. The GLW dataset is well suited to this gap because it preserves repeated molecular measurements, visit gaps, interval survival structure, and privacy-protected patient identifiers.

The methodological contribution is therefore to analyze DMR as a longitudinal-event process rather than as a single exact response date. Joint longitudinal-survival models and related mixed-observation methods provide a natural framework for connecting biomarker trajectories to event risk \cite{Wulfsohn1997JointModel,Rizopoulos2016JMbayes,Lovblom2026MixedObservationJointModel,Niu2026QuantileJointModel}. In this setting, the GLW study can complement guideline and TFR literature by focusing on the real-world data structure generated by clinical monitoring: sparse and irregular visits, interval-observed response times, molecular assay limits, and partially observed clinical covariates.

